% !TEX program = xelatex
\documentclass[10pt,a4paper]{article}
\usepackage[utf8]{inputenc}
\usepackage[T1]{fontenc}
\usepackage[french]{babel}
\usepackage[left=0.6cm,right=0.6cm,top=1cm,bottom=1cm]{geometry}
\usepackage{amsmath}
\usepackage{amssymb}
\usepackage{array}
\usepackage{graphicx}
\usepackage{xcolor}
\usepackage{titlesec}
\usepackage{enumitem}
\usepackage{multicol}
\usepackage{pifont}
% Badges « Exercice n » en cursive Playwrite FR Moderne (police du site),
% chargée à 85 % de la taille courante (XeLaTeX ou LuaLaTeX ; gras sous pdfLaTeX).
\makeatletter
\ifdefined\XeTeXrevision\def\ExoFont@ok{}\fi
\ifdefined\directlua\def\ExoFont@ok{}\fi
\ifdefined\ExoFont@ok
\newcommand{\ExoFont}{\font\ExoFont@f="[PlaywriteFRModerne.ttf]" at \dimexpr\f@size pt*85/100\relax\ExoFont@f}
\else
\newcommand{\ExoFont}{\bfseries}
\fi
\makeatother
% \ExoTexte{...} : texte du badge en cursive, avec la hauteur d'une majuscule de la
% police courante (la cursive est plus haute et ferait grandir chaque badge).
\newcommand{\ExoTexte}[1]{\raisebox{0pt}[\fontcharht\font`E][0pt]{\ExoFont #1}}
\usepackage{frcursive}
% Correctif du bug de crénage de frcursive sous XeLaTeX (repris de CLT2627.sty) :
% sans lui, certaines paires de lettres consécutives dans \cursive explosent
% (une lettre projetée loin à droite), ce qui coupe des mots comme "Représenter".
\makeatletter
\ifdefined\XeTeXcharclass
\XeTeXcharclass `a 201 \XeTeXcharclass `b 201 \XeTeXcharclass `c 201
\XeTeXcharclass `d 201 \XeTeXcharclass `e 201 \XeTeXcharclass `f 201
\XeTeXcharclass `g 201 \XeTeXcharclass `h 201 \XeTeXcharclass `i 201
\XeTeXcharclass `j 201 \XeTeXcharclass `k 201 \XeTeXcharclass `l 201
\XeTeXcharclass `m 201 \XeTeXcharclass `n 201 \XeTeXcharclass `o 201
\XeTeXcharclass `p 201 \XeTeXcharclass `q 201 \XeTeXcharclass `r 201
\XeTeXcharclass `s 201 \XeTeXcharclass `t 201 \XeTeXcharclass `u 201
\XeTeXcharclass `v 201 \XeTeXcharclass `w 201 \XeTeXcharclass `x 201
\XeTeXcharclass `y 201 \XeTeXcharclass `z 201
\XeTeXcharclass `A 201 \XeTeXcharclass `B 201 \XeTeXcharclass `C 201
\XeTeXcharclass `D 201 \XeTeXcharclass `E 201 \XeTeXcharclass `F 201
\XeTeXcharclass `G 201 \XeTeXcharclass `H 201 \XeTeXcharclass `I 201
\XeTeXcharclass `J 201 \XeTeXcharclass `K 201 \XeTeXcharclass `L 201
\XeTeXcharclass `M 201 \XeTeXcharclass `N 201 \XeTeXcharclass `O 201
\XeTeXcharclass `P 201 \XeTeXcharclass `Q 201 \XeTeXcharclass `R 201
\XeTeXcharclass `S 201 \XeTeXcharclass `T 201 \XeTeXcharclass `U 201
\XeTeXcharclass `V 201 \XeTeXcharclass `W 201 \XeTeXcharclass `X 201
\XeTeXcharclass `Y 201 \XeTeXcharclass `Z 201
\XeTeXcharclass `\é 201 \XeTeXcharclass `\è 201 \XeTeXcharclass `\ê 201
\XeTeXcharclass `\ë 201 \XeTeXcharclass `\É 201 \XeTeXcharclass `\È 201
\XeTeXcharclass `\Ê 201 \XeTeXcharclass `\Ë 201 \XeTeXcharclass `\à 201
\XeTeXcharclass `\â 201 \XeTeXcharclass `\ä 201 \XeTeXcharclass `\À 201
\XeTeXcharclass `\Â 201 \XeTeXcharclass `\Ä 201 \XeTeXcharclass `\ô 201
\XeTeXcharclass `\ö 201 \XeTeXcharclass `\Ô 201 \XeTeXcharclass `\Ö 201
\XeTeXcharclass `\î 201 \XeTeXcharclass `\ï 201 \XeTeXcharclass `\Î 201
\XeTeXcharclass `\Ï 201 \XeTeXcharclass `\ù 201 \XeTeXcharclass `\û 201
\XeTeXcharclass `\ü 201 \XeTeXcharclass `\Ù 201 \XeTeXcharclass `\Û 201
\XeTeXcharclass `\Ü 201 \XeTeXcharclass `\ç 201 \XeTeXcharclass `\Ç 201
\XeTeXcharclass `\œ 201 \XeTeXcharclass `\Œ 201 \XeTeXcharclass `\æ 201
\XeTeXcharclass `\Æ 201
\let\cursive@orig\cursive
\renewcommand{\cursive}{%
  \cursive@orig
  \XeTeXinterchartoks 201 201 = {\kern0pt }%
}
\fi
\makeatother
\usepackage{tikz}
\usetikzlibrary{calc,arrows.meta,shapes.misc}
\tikzset{ptcross/.style={cross out,draw,thick,minimum size=3pt,inner sep=0pt}}
\usepackage{tcolorbox}
\tcbuselibrary{skins,breakable}
\usepackage[hidelinks]{hyperref}
\usepackage{fontawesome5}
\definecolor{calcblue}{HTML}{1C7ED6}
\definecolor{diffgreen}{HTML}{40C057}
\definecolor{diffblue}{HTML}{228BE6}
\definecolor{diffred}{HTML}{FA5252}
\definecolor{skyblue}{HTML}{339AF0}
\definecolor{objorange}{HTML}{F76707}
\newcommand{\starfull}[1]{{\color{#1}\ding{72}}}
\newcommand{\starempty}{{\color{gray!45}\ding{73}}}
\newcommand{\Vj}[1]{\overrightarrow{\text{#1}}}
\newcommand{\VBlank}{\rule[0.15ex]{0.55cm}{0.35pt}}
\newcommand{\Pt}[1]{\textsf{\footnotesize #1}}
\renewcommand{\familydefault}{\sfdefault}
\DeclareSymbolFont{sfnumbers}{OT1}{cmss}{m}{n}
\SetSymbolFont{sfnumbers}{bold}{OT1}{cmss}{bx}{n}
\DeclareMathSymbol{0}{\mathalpha}{sfnumbers}{`0}
\DeclareMathSymbol{1}{\mathalpha}{sfnumbers}{`1}
\DeclareMathSymbol{2}{\mathalpha}{sfnumbers}{`2}
\DeclareMathSymbol{3}{\mathalpha}{sfnumbers}{`3}
\DeclareMathSymbol{4}{\mathalpha}{sfnumbers}{`4}
\DeclareMathSymbol{5}{\mathalpha}{sfnumbers}{`5}
\DeclareMathSymbol{6}{\mathalpha}{sfnumbers}{`6}
\DeclareMathSymbol{7}{\mathalpha}{sfnumbers}{`7}
\DeclareMathSymbol{8}{\mathalpha}{sfnumbers}{`8}
\DeclareMathSymbol{9}{\mathalpha}{sfnumbers}{`9}
\setlength{\parindent}{0pt}
\graphicspath{{./}}
\newcommand*{\textcircledd}[1]{\begin{tikzpicture}[baseline=(letter.base)]\node[draw,circle,inner sep=1pt] (letter) {#1};\end{tikzpicture}}
\newcommand{\numeroCerne}[1]{\protect\textcircledd{\bfseries\scriptsize #1}}
\setlist[enumerate,1]{label=\numeroCerne{\arabic*}, leftmargin=*, labelsep=6pt, itemsep=3pt, topsep=2pt, partopsep=0pt, parsep=0pt}

\NewDocumentEnvironment{Exo}{O{breakable} m m m}{%
  \begin{tcolorbox}[
    enhanced,
    #1,
    colback=white,
    colframe=white,
    boxrule=0pt,
    arc=4pt, outer arc=4pt, sharp corners=south,
    left=8pt, right=8pt, top=5pt, bottom=3pt,
    before skip=4pt, after skip=4pt,
  ]
  \noindent\tcbox[on line, colback=skyblue, colframe=skyblue, boxsep=1pt, left=4pt, right=4pt, top=1pt, bottom=1pt, arc=2pt]{\textcolor{white}{\ExoTexte{Exercice~#2}}}~#3\hfill\mbox{\color{skyblue}\normalsize\cursive\hyphenchar\font=-1 #4}\par\vspace{3pt}%
}{%
  \end{tcolorbox}
}

\NewDocumentEnvironment{ExoCont}{O{unbreakable}}{%
  \begin{tcolorbox}[
    enhanced,
    #1,
    colback=white,
    colframe=white,
    boxrule=0pt,
    arc=4pt, outer arc=4pt, sharp corners=south,
    left=8pt, right=8pt, top=5pt, bottom=3pt,
    before skip=4pt, after skip=4pt,
  ]
}{%
  \end{tcolorbox}
}

\newcommand{\Objectif}[2]{%
\vspace{4pt}
\noindent{\large\color{objorange}\bfseries Objectif #1~:}\\{\large\color{objorange}#2}\par
\vspace{2pt}
}

\NewDocumentCommand{\FeuilletTitre}{m}{%
  \begin{center}
  \begin{tcolorbox}[
    enhanced,
    width=0.58\linewidth,
    colback=white,
    colframe=black,
    boxrule=0.8pt,
    arc=8pt,
    sharp corners=downhill,
    left=10pt, right=10pt, top=8pt, bottom=6pt,
    halign=center,
    overlay={
      \node[fill=white, anchor=west, font=\normalsize, inner xsep=4pt]
        at ([xshift=14pt]frame.north west) {Feuille d'entra\^inement};
    },
  ]
  {\huge #1}
  \end{tcolorbox}
  \end{center}
}

\pagestyle{empty}
\definecolor{repred}{HTML}{E03131}
\newcommand{\Rep}[1]{{\color{repred}#1}}
% Ligne de corrigé : numéro d'exercice en pastille bleu ciel (comme le corrigé de la feuille 3).
\newcommand{\Corr}[1]{%
\par\vspace{8pt}
\noindent\tcbox[on line, colback=skyblue, colframe=skyblue, boxsep=1pt, left=4pt, right=4pt, top=1pt, bottom=1pt, arc=2pt]{\textcolor{white}{\ExoTexte{Exercice~#1}}}\par\nobreak\vspace{3pt}\nobreak
}
\begin{document}

\FeuilletTitre{Vecteurs colin\'eaires}

\Objectif{5}{Produit d'un vecteur par un nombre r\'eel et colin\'earit\'e de deux vecteurs}

\begin{Exo}[breakable]{1}{\starfull{diffblue}\starfull{diffblue}\starempty}{Repr\'esenter}
\renewcommand{\Pt}[1]{\textsf{#1}}
Exprimer le vecteur $\Vj{AC}$ en fonction du vecteur $\Vj{AB}$ dans chacun des cas suivants~:

\medskip

\noindent\begin{minipage}[t]{0.48\linewidth}\vspace{0pt}
\noindent\numeroCerne{1}\par
{\centering
\begin{tikzpicture}[x=1.286cm,y=0.5cm]
\draw (0,0) -- (7,0);
\foreach \x in {2,4,5} \draw (\x,-0.08) -- (\x,0.08);
\node[ptcross] at (3,0) {}; \node[above] at (3,0) {\Pt{A}};
\node[ptcross] at (6,0) {}; \node[above] at (6,0) {\Pt{B}};
\node[ptcross] at (1,0) {}; \node[above] at (1,0) {\Pt{C}};
\end{tikzpicture}\par}
\end{minipage}\hfill
\begin{minipage}[t]{0.48\linewidth}\vspace{0pt}
\noindent\numeroCerne{2}\par
{\centering
\begin{tikzpicture}[x=1.286cm,y=0.5cm]
\draw (0,0) -- (7,0);
\foreach \x in {2,4,5} \draw (\x,-0.08) -- (\x,0.08);
\node[ptcross] at (1,0) {}; \node[above] at (1,0) {\Pt{A}};
\node[ptcross] at (3,0) {}; \node[above] at (3,0) {\Pt{B}};
\node[ptcross] at (6,0) {}; \node[above] at (6,0) {\Pt{C}};
\end{tikzpicture}\par}
\end{minipage}

\vspace{1.1cm}

\noindent\begin{minipage}[t]{0.48\linewidth}\vspace{0pt}
\noindent\numeroCerne{3}\par
{\centering
\begin{tikzpicture}[x=1.125cm,y=0.5cm]
\draw (0,0) -- (8,0);
\foreach \x in {1,3,5,6} \draw (\x,-0.08) -- (\x,0.08);
\node[ptcross] at (7,0) {}; \node[above] at (7,0) {\Pt{A}};
\node[ptcross] at (2,0) {}; \node[above] at (2,0) {\Pt{B}};
\node[ptcross] at (4,0) {}; \node[above] at (4,0) {\Pt{C}};
\end{tikzpicture}\par}
\end{minipage}\hfill
\begin{minipage}[t]{0.48\linewidth}\vspace{0pt}
\noindent\numeroCerne{4}\par
{\centering
\begin{tikzpicture}[x=1.500cm,y=0.5cm]
\draw (-1,0) -- (5,0);
\foreach \x in {1,4} \draw (\x,-0.08) -- (\x,0.08);
\node[ptcross] at (2,0) {}; \node[above] at (2,0) {\Pt{A}};
\node[ptcross] at (3,0) {}; \node[below] at (3,0) {\Pt{B}};
\node[ptcross] at (0,0) {}; \node[above] at (0,0) {\Pt{C}};
\end{tikzpicture}\par}
\end{minipage}

\vspace{1.1cm}

\noindent\begin{minipage}[t]{0.48\linewidth}\vspace{0pt}
\noindent\numeroCerne{5}\par
{\centering
\begin{tikzpicture}[x=0.750cm,y=0.5cm]
\draw (-4,0) -- (8,0);
\foreach \x in {-2,-1,0,1,2,4,5,6} \draw (\x,-0.08) -- (\x,0.08);
\node[ptcross] at (3,0) {}; \node[above] at (3,0) {\Pt{A}};
\node[ptcross] at (7,0) {}; \node[above] at (7,0) {\Pt{B}};
\node[ptcross] at (-3,0) {}; \node[above] at (-3,0) {\Pt{C}};
\end{tikzpicture}\par}
\end{minipage}\hfill
\begin{minipage}[t]{0.48\linewidth}\vspace{0pt}
\noindent\numeroCerne{6}\par
{\centering
\begin{tikzpicture}[x=1.125cm,y=0.5cm]
\draw (0,0) -- (8,0);
\foreach \x in {2,3,4,6} \draw (\x,-0.08) -- (\x,0.08);
\node[ptcross] at (1,0) {}; \node[above] at (1,0) {\Pt{A}};
\node[ptcross] at (7,0) {}; \node[above] at (7,0) {\Pt{B}};
\node[ptcross] at (5,0) {}; \node[above] at (5,0) {\Pt{C}};
\end{tikzpicture}\par}
\end{minipage}

\vspace{1.1cm}

\noindent\begin{minipage}[t]{0.48\linewidth}\vspace{0pt}
\noindent\numeroCerne{7}\par
{\centering
\begin{tikzpicture}[x=1.125cm,y=0.5cm]
\draw (-2,0) -- (6,0);
\foreach \x in {0,1,2,4} \draw (\x,-0.08) -- (\x,0.08);
\node[ptcross] at (5,0) {}; \node[above] at (5,0) {\Pt{A}};
\node[ptcross] at (3,0) {}; \node[above] at (3,0) {\Pt{B}};
\node[ptcross] at (-1,0) {}; \node[above] at (-1,0) {\Pt{C}};
\end{tikzpicture}\par}
\end{minipage}\hfill
\begin{minipage}[t]{0.48\linewidth}\vspace{0pt}
\noindent\numeroCerne{8}\par
{\centering
\begin{tikzpicture}[x=1.286cm,y=0.5cm]
\draw (2,0) -- (9,0);
\foreach \x in {5,6,7} \draw (\x,-0.08) -- (\x,0.08);
\node[ptcross] at (4,0) {}; \node[above] at (4,0) {\Pt{A}};
\node[ptcross] at (8,0) {}; \node[above] at (8,0) {\Pt{B}};
\node[ptcross] at (3,0) {}; \node[below] at (3,0) {\Pt{C}};
\end{tikzpicture}\par}
\end{minipage}

\end{Exo}

\begin{Exo}[breakable]{2}{\starfull{diffblue}\starfull{diffblue}\starempty}{Repr\'esenter}
\renewcommand{\Pt}[1]{\textsf{#1}}
Dans chacun des cas suivants, repr\'esenter le vecteur $\Vj{AC}$ tel que~:

\medskip

\noindent\begin{minipage}[t]{0.48\linewidth}\vspace{0pt}
\noindent\numeroCerne{1}~$\Vj{AC}=\dfrac{3}{2}\Vj{AB}$\par
{\centering
\begin{tikzpicture}[x=1.286cm,y=0.5cm]
\draw (0,0) -- (7,0);
\foreach \x in {1,3,5,6} \draw (\x,-0.08) -- (\x,0.08);
\node[ptcross] at (2,0) {}; \node[above] at (2,0) {\Pt{A}};
\node[ptcross] at (4,0) {}; \node[above] at (4,0) {\Pt{B}};
\end{tikzpicture}\par}
\end{minipage}\hfill
\begin{minipage}[t]{0.48\linewidth}\vspace{0pt}
\noindent\numeroCerne{2}~$\Vj{AC}=-\dfrac{2}{3}\Vj{AB}$\par
{\centering
\begin{tikzpicture}[x=1.125cm,y=0.5cm]
\draw (0,0) -- (8,0);
\foreach \x in {1,2,3,5,6} \draw (\x,-0.08) -- (\x,0.08);
\node[ptcross] at (4,0) {}; \node[above] at (4,0) {\Pt{A}};
\node[ptcross] at (7,0) {}; \node[above] at (7,0) {\Pt{B}};
\end{tikzpicture}\par}
\end{minipage}

\vspace{1.1cm}

\noindent\begin{minipage}[t]{0.48\linewidth}\vspace{0pt}
\noindent\numeroCerne{3}~$\Vj{AC}=\dfrac{5}{4}\Vj{AB}$\par
{\centering
\begin{tikzpicture}[x=1.125cm,y=0.5cm]
\draw (0,0) -- (8,0);
\foreach \x in {2,3,4,6,7} \draw (\x,-0.08) -- (\x,0.08);
\node[ptcross] at (1,0) {}; \node[above] at (1,0) {\Pt{A}};
\node[ptcross] at (5,0) {}; \node[above] at (5,0) {\Pt{B}};
\end{tikzpicture}\par}
\end{minipage}\hfill
\begin{minipage}[t]{0.48\linewidth}\vspace{0pt}
\noindent\numeroCerne{4}~$\Vj{AC}=-\dfrac{1}{3}\Vj{AB}$\par
{\centering
\begin{tikzpicture}[x=0.900cm,y=0.5cm]
\draw (0,0) -- (10,0);
\foreach \x in {2,3,4,5,6,8,9} \draw (\x,-0.08) -- (\x,0.08);
\node[ptcross] at (7,0) {}; \node[above] at (7,0) {\Pt{A}};
\node[ptcross] at (1,0) {}; \node[above] at (1,0) {\Pt{B}};
\end{tikzpicture}\par}
\end{minipage}

\vspace{1.1cm}

\noindent\begin{minipage}[t]{0.48\linewidth}\vspace{0pt}
\noindent\numeroCerne{5}~$\Vj{AC}=\dfrac{3}{5}\Vj{AB}$\par
{\centering
\begin{tikzpicture}[x=1.000cm,y=0.5cm]
\draw (0,0) -- (9,0);
\foreach \x in {1,3,4,5,6,8} \draw (\x,-0.08) -- (\x,0.08);
\node[ptcross] at (7,0) {}; \node[above] at (7,0) {\Pt{A}};
\node[ptcross] at (2,0) {}; \node[above] at (2,0) {\Pt{B}};
\end{tikzpicture}\par}
\end{minipage}\hfill
\begin{minipage}[t]{0.48\linewidth}\vspace{0pt}
\noindent\numeroCerne{6}~$\Vj{AC}=-\dfrac{5}{2}\Vj{AB}$\par
{\centering
\begin{tikzpicture}[x=1.000cm,y=0.5cm]
\draw (0,0) -- (9,0);
\foreach \x in {1,2,3,4,5,7} \draw (\x,-0.08) -- (\x,0.08);
\node[ptcross] at (6,0) {}; \node[above] at (6,0) {\Pt{A}};
\node[ptcross] at (8,0) {}; \node[above] at (8,0) {\Pt{B}};
\end{tikzpicture}\par}
\end{minipage}

\vspace{1.1cm}

\noindent\begin{minipage}[t]{0.48\linewidth}\vspace{0pt}
\noindent\numeroCerne{7}~$\Vj{AC}=-3\Vj{AB}$\par
{\centering
\begin{tikzpicture}[x=0.900cm,y=0.5cm]
\draw (0,0) -- (10,0);
\foreach \x in {2,4,5,6,7,8,9} \draw (\x,-0.08) -- (\x,0.08);
\node[ptcross] at (3,0) {}; \node[above] at (3,0) {\Pt{A}};
\node[ptcross] at (1,0) {}; \node[above] at (1,0) {\Pt{B}};
\end{tikzpicture}\par}
\end{minipage}\hfill
\begin{minipage}[t]{0.48\linewidth}\vspace{0pt}
\noindent\numeroCerne{8}~$\Vj{AC}=\dfrac{7}{4}\Vj{AB}$\par
{\centering
\begin{tikzpicture}[x=0.900cm,y=0.5cm]
\draw (0,0) -- (10,0);
\foreach \x in {2,3,4,6,7,8,9} \draw (\x,-0.08) -- (\x,0.08);
\node[ptcross] at (1,0) {}; \node[above] at (1,0) {\Pt{A}};
\node[ptcross] at (5,0) {}; \node[above] at (5,0) {\Pt{B}};
\end{tikzpicture}\par}
\end{minipage}

\end{Exo}

\newpage

\begin{center}
{\Large Corrig\'e}
\end{center}

\renewcommand{\Pt}[1]{\textsf{#1}}
\Corr{1}
On compte les longueurs en graduations. Le signe est n\'egatif quand $\Vj{AC}$ et $\Vj{AB}$ sont de sens oppos\'e.
\begin{multicols}{2}\raggedright
\begin{enumerate}[itemsep=9pt]
\item \begin{minipage}[t]{\linewidth}Sens oppos\'e, $AC=2$ et $AB=3$.\\[2pt]
$\Vj{AC}=-\dfrac{2}{3}\,\Vj{AB}$\end{minipage}
\item \begin{minipage}[t]{\linewidth}M\^eme sens, $AC=5$ et $AB=2$.\\[2pt]
$\Vj{AC}=\dfrac{5}{2}\,\Vj{AB}$\end{minipage}
\item \begin{minipage}[t]{\linewidth}M\^eme sens, $AC=3$ et $AB=5$.\\[2pt]
$\Vj{AC}=\dfrac{3}{5}\,\Vj{AB}$\end{minipage}
\item \begin{minipage}[t]{\linewidth}Sens oppos\'e, $AC=2$ et $AB=1$.\\[2pt]
$\Vj{AC}=-2\,\Vj{AB}$\end{minipage}
\item \begin{minipage}[t]{\linewidth}Sens oppos\'e, $AC=6$ et $AB=4$.\\[2pt]
$\Vj{AC}=-\dfrac{6}{4}\,\Vj{AB}=-\dfrac{3}{2}\,\Vj{AB}$\end{minipage}
\item \begin{minipage}[t]{\linewidth}M\^eme sens, $AC=4$ et $AB=6$.\\[2pt]
$\Vj{AC}=\dfrac{4}{6}\,\Vj{AB}=\dfrac{2}{3}\,\Vj{AB}$\end{minipage}
\item \begin{minipage}[t]{\linewidth}M\^eme sens, $AC=6$ et $AB=2$.\\[2pt]
$\Vj{AC}=\dfrac{6}{2}\,\Vj{AB}=3\,\Vj{AB}$\end{minipage}
\item \begin{minipage}[t]{\linewidth}Sens oppos\'e, $AC=1$ et $AB=4$.\\[2pt]
$\Vj{AC}=-\dfrac{1}{4}\,\Vj{AB}$\end{minipage}
\end{enumerate}
\end{multicols}

\Corr{2}
On place $C$ sur la droite $(AB)$ : du c\^ot\'e de $B$ si le nombre est positif, \`a l'oppos\'e sinon.
\begin{multicols}{2}\raggedright
\begin{enumerate}[itemsep=9pt]
\item \begin{minipage}[t]{\linewidth}\raggedright $\Vj{AC}=\dfrac{3}{2}\Vj{AB}$~: $\dfrac{3}{2}>0$ donc m\^eme sens.\\
$AB=2$, donc $AC=\dfrac{3}{2}\times2=3$ graduations.\par\vspace{2pt}
{\centering
\begin{tikzpicture}[x=1.171cm,y=0.5cm]
\draw (0,0) -- (7,0);
\foreach \x in {1,3,6} \draw (\x,-0.08) -- (\x,0.08);
\node[ptcross] at (2,0) {}; \node[above] at (2,0) {\Pt{A}};
\node[ptcross] at (4,0) {}; \node[above] at (4,0) {\Pt{B}};
\node[ptcross,repred] at (5,0) {}; \node[below,repred] at (5,0) {\Pt{C}};
\draw[-{Stealth[length=6pt]},thick,repred] (2,1.7) -- (5,1.7);
\end{tikzpicture}\par}\end{minipage}
\item \begin{minipage}[t]{\linewidth}\raggedright $\Vj{AC}=-\dfrac{2}{3}\Vj{AB}$~: $-\dfrac{2}{3}<0$ donc sens oppos\'e.\\
$AB=3$, donc $AC=\dfrac{2}{3}\times3=2$ graduations.\par\vspace{2pt}
{\centering
\begin{tikzpicture}[x=1.025cm,y=0.5cm]
\draw (0,0) -- (8,0);
\foreach \x in {1,3,5,6} \draw (\x,-0.08) -- (\x,0.08);
\node[ptcross] at (4,0) {}; \node[above] at (4,0) {\Pt{A}};
\node[ptcross] at (7,0) {}; \node[above] at (7,0) {\Pt{B}};
\node[ptcross,repred] at (2,0) {}; \node[above,repred] at (2,0) {\Pt{C}};
\draw[-{Stealth[length=6pt]},thick,repred] (4,1.7) -- (2,1.7);
\end{tikzpicture}\par}\end{minipage}
\item \begin{minipage}[t]{\linewidth}\raggedright $\Vj{AC}=\dfrac{5}{4}\Vj{AB}$~: $\dfrac{5}{4}>0$ donc m\^eme sens.\\
$AB=4$, donc $AC=\dfrac{5}{4}\times4=5$ graduations.\par\vspace{2pt}
{\centering
\begin{tikzpicture}[x=1.025cm,y=0.5cm]
\draw (0,0) -- (8,0);
\foreach \x in {2,3,4,7} \draw (\x,-0.08) -- (\x,0.08);
\node[ptcross] at (1,0) {}; \node[above] at (1,0) {\Pt{A}};
\node[ptcross] at (5,0) {}; \node[above] at (5,0) {\Pt{B}};
\node[ptcross,repred] at (6,0) {}; \node[below,repred] at (6,0) {\Pt{C}};
\draw[-{Stealth[length=6pt]},thick,repred] (1,1.7) -- (6,1.7);
\end{tikzpicture}\par}\end{minipage}
\item \begin{minipage}[t]{\linewidth}\raggedright $\Vj{AC}=-\dfrac{1}{3}\Vj{AB}$~: $-\dfrac{1}{3}<0$ donc sens oppos\'e.\\
$AB=6$, donc $AC=\dfrac{1}{3}\times6=2$ graduations.\par\vspace{2pt}
{\centering
\begin{tikzpicture}[x=0.820cm,y=0.5cm]
\draw (0,0) -- (10,0);
\foreach \x in {2,3,4,5,6,8} \draw (\x,-0.08) -- (\x,0.08);
\node[ptcross] at (7,0) {}; \node[above] at (7,0) {\Pt{A}};
\node[ptcross] at (1,0) {}; \node[above] at (1,0) {\Pt{B}};
\node[ptcross,repred] at (9,0) {}; \node[above,repred] at (9,0) {\Pt{C}};
\draw[-{Stealth[length=6pt]},thick,repred] (7,1.7) -- (9,1.7);
\end{tikzpicture}\par}\end{minipage}
\item \begin{minipage}[t]{\linewidth}\raggedright $\Vj{AC}=\dfrac{3}{5}\Vj{AB}$~: $\dfrac{3}{5}>0$ donc m\^eme sens.\\
$AB=5$, donc $AC=\dfrac{3}{5}\times5=3$ graduations.\par\vspace{2pt}
{\centering
\begin{tikzpicture}[x=0.911cm,y=0.5cm]
\draw (0,0) -- (9,0);
\foreach \x in {1,3,5,6,8} \draw (\x,-0.08) -- (\x,0.08);
\node[ptcross] at (7,0) {}; \node[above] at (7,0) {\Pt{A}};
\node[ptcross] at (2,0) {}; \node[above] at (2,0) {\Pt{B}};
\node[ptcross,repred] at (4,0) {}; \node[above,repred] at (4,0) {\Pt{C}};
\draw[-{Stealth[length=6pt]},thick,repred] (7,1.7) -- (4,1.7);
\end{tikzpicture}\par}\end{minipage}
\item \begin{minipage}[t]{\linewidth}\raggedright $\Vj{AC}=-\dfrac{5}{2}\Vj{AB}$~: $-\dfrac{5}{2}<0$ donc sens oppos\'e.\\
$AB=2$, donc $AC=\dfrac{5}{2}\times2=5$ graduations.\par\vspace{2pt}
{\centering
\begin{tikzpicture}[x=0.911cm,y=0.5cm]
\draw (0,0) -- (9,0);
\foreach \x in {2,3,4,5,7} \draw (\x,-0.08) -- (\x,0.08);
\node[ptcross] at (6,0) {}; \node[above] at (6,0) {\Pt{A}};
\node[ptcross] at (8,0) {}; \node[above] at (8,0) {\Pt{B}};
\node[ptcross,repred] at (1,0) {}; \node[above,repred] at (1,0) {\Pt{C}};
\draw[-{Stealth[length=6pt]},thick,repred] (6,1.7) -- (1,1.7);
\end{tikzpicture}\par}\end{minipage}
\item \begin{minipage}[t]{\linewidth}\raggedright $\Vj{AC}=-3\Vj{AB}$~: $-3<0$ donc sens oppos\'e.\\
$AB=2$, donc $AC=3\times2=6$ graduations.\par\vspace{2pt}
{\centering
\begin{tikzpicture}[x=0.820cm,y=0.5cm]
\draw (0,0) -- (10,0);
\foreach \x in {2,4,5,6,7,8} \draw (\x,-0.08) -- (\x,0.08);
\node[ptcross] at (3,0) {}; \node[above] at (3,0) {\Pt{A}};
\node[ptcross] at (1,0) {}; \node[above] at (1,0) {\Pt{B}};
\node[ptcross,repred] at (9,0) {}; \node[above,repred] at (9,0) {\Pt{C}};
\draw[-{Stealth[length=6pt]},thick,repred] (3,1.7) -- (9,1.7);
\end{tikzpicture}\par}\end{minipage}
\item \begin{minipage}[t]{\linewidth}\raggedright $\Vj{AC}=\dfrac{7}{4}\Vj{AB}$~: $\dfrac{7}{4}>0$ donc m\^eme sens.\\
$AB=4$, donc $AC=\dfrac{7}{4}\times4=7$ graduations.\par\vspace{2pt}
{\centering
\begin{tikzpicture}[x=0.820cm,y=0.5cm]
\draw (0,0) -- (10,0);
\foreach \x in {2,3,4,6,7,9} \draw (\x,-0.08) -- (\x,0.08);
\node[ptcross] at (1,0) {}; \node[above] at (1,0) {\Pt{A}};
\node[ptcross] at (5,0) {}; \node[above] at (5,0) {\Pt{B}};
\node[ptcross,repred] at (8,0) {}; \node[above,repred] at (8,0) {\Pt{C}};
\draw[-{Stealth[length=6pt]},thick,repred] (1,1.7) -- (8,1.7);
\end{tikzpicture}\par}\end{minipage}
\end{enumerate}
\end{multicols}

\end{document}
